\section{Related work}\label{sec:related}

\paragraph{Batteries for demand-charge management} Sizing and operating storage for peak shaving is a well-studied problem~\cite{oudalov2007sizing,leadbetter2012battery,papadopoulos2020peak}. Schedules are typically computed from a load (and, with photovoltaics, a generation) forecast by linear or mixed-integer programming, so that forecast errors reduce the realised saving~\cite{hanna2014dispatch}. Recent studies optimise storage and vehicle-to-building resources for commercial buildings, in one case with machine-learned load forecasts~\cite{ghafoori2023v2b,jiang2026v2b}, and handle forecast error with robust model predictive control~\cite{mary2025rmpc}. These studies optimise a controller for one forecaster or a few buildings; they do not ask how the choice among forecasters, judged by accuracy, relates to the bill across a population of buildings.

\paragraph{Evaluating load forecasts} Building and smart-meter load forecasting is evaluated almost entirely by error metrics~\cite{hong2016probabilistic,wang2019smartmeter}, including in the largest recent benchmark~\cite{emami2023buildingsbench}. Soman et al.~\cite{soman2020peak} distinguish predicting which hours will be the peak from predicting the whole load curve, and argue that the former is what battery-based optimisation on a campus microgrid needs. Gneiting~\cite{gneiting2011point} shows more generally that a point forecast is optimal only relative to a scoring function, so that ranking forecasts by MAE implicitly assumes that the decision maker's loss is symmetric and linear in the error.

\paragraph{Decision-focused and value-oriented forecasting} Decision-focused learning trains a predictor on the cost of the downstream optimisation instead of on its error~\cite{elmachtoub2022spo,mandi2024dfl,sadana2025survey}. In energy, value-oriented forecasting has been applied to trading renewable output~\cite{stratigakos2021value}, to demand-charge storage with a jointly trained forecaster and controller~\cite{kim2025neural}, and to residential photovoltaic-battery scheduling on twenty buildings, where it lowered costs by a few per cent at a loss of accuracy~\cite{depoortere2026dfl}. The present study is complementary. It proposes no new training method; it measures how well the accuracy of conventional forecasters predicts their decision value on \res{nKept} buildings and decomposes the gap by type of error and by controller, which indicates when decision-focused training, or simply a different choice among existing forecasters, is worth the effort.
